%!TEX root = hems_proj.tex
%%%%%%%%%%%%%%%%%%%%%%%%%%%%%%%%%%%%%%%%%%%%%%%%%%%%%%%%%%%%%%%%%%%%%%%
\chapter{Bevezetés és problémafelvetés}\label{ch:BEVEZETES}
%%%%%%%%%%%%%%%%%%%%%%%%%%%%%%%%%%%%%%%%%%%%%%%%%%%%%%%%%%%%%%%%%%%%%%%

\begin{osszefoglal}
A fejezet bemutatja a dolgozat tárgyát, a keresletoldali szabályozás jelentőségét, valamint a kitűzött kutatási célokat.
\end{osszefoglal}

A globális energetikai átalakulás, a decentralizált termelés térnyerése és a fenntarthatósági elvárások a Smart Grid technológiák felé terelik a villamosenergia-hálózatokat. Ebben a környezetben az otthoni energiamenedzsment-rendszer (Home Energy Management System, HEMS) a lakossági oldal egyik fontos eszköze, mivel összehangolja a fogyasztókat, az energiatárolást és a megújuló termelést, miközben kétirányú energiaáramlást és kommunikációt tesz lehetővé. A dolgozat középpontjában nem maga a HEMS hardvere áll, hanem az a kérdés, hogy a benne futó optimalizáló algoritmus milyen bemeneti körülmények mellett hoz megbízható döntéseket.

A dolgozat elsődleges célja egy olyan optimalizálási modell kidolgozása és értékelése, amely a háztartási eszközök összehangolt vezérlésével és a megújuló források jobb kihasználásával csökkenti a napi villamosenergia-költséget, miközben tiszteletben tartja a fizikai és hálózati korlátokat. A megközelítés emellett az energiatárolók intelligens irányítását és a hálózati terhelési csúcsok mérséklését is szem előtt tartja, a felhasználói komfort fenntartása mellett.

A keresletoldali szabályozás (Demand Side Management, DSM) és különösen a terheléseltolás (Load Shifting) több, egymást erősítő előnnyel jár \cite{Fadaee2012}. A legközvetlenebb hozadék a számottevő, akár 20--40\%-os költségmegtakarítás, ám ezen túl a módszer hozzájárul a hálózati csúcsok csökkentéséhez, szélesebb körben alkalmazva pedig mérsékli a fosszilis erőművek igénybevételét, ezáltal a károsanyag-kibocsátást is.

A HEMS-optimalizálás a gyakorlatban több, egymással versengő cél kezelését jelenti: a költség minimalizálását, a felhasználói komfort fenntartását és a hálózati terhelés csökkentését. A lakossági fogyasztás időben erősen ingadozó, a napelemes termelés pedig időjárásfüggő és bizonytalan, ezért a feladat nem egyszerű ütemezési probléma, hanem korlátokkal terhelt, vegyes (diszkrét és folytonos) döntési változókat tartalmazó optimalizációs modellként kezelhető \cite{Ismail2014, Fadaee2012}. A téma aktualitását a megújuló kapacitások bővülése és a dinamikus tarifák terjedése adja, amelyek a fogyasztói oldalon valódi rugalmasságot tesznek lehetővé.

Mindezek alapján a dolgozat egy olyan optimalizációs keretet mutat be, amely természetből inspirált metaheurisztikák, nevezetesen a genetikus algoritmus (GA), a részecskeraj-optimalizálás (PSO), a szürkefarkas-optimalizáló (GWO) és egy hibrid GA--PSO eljárás segítségével csökkenti a napi villamosenergia-költséget. A munka legfontosabb hozzájárulása nem önmagában az algoritmusok alkalmazása, hanem az, hogy a teljesítményüket valós és szintetikus bemenetek szisztematikus összehasonlításában vizsgálja, így a hangsúly az algoritmusok robusztusságán és gyakorlati általánosíthatóságán van \cite{Lazzari2023}.

\paragraph{A mesterséges intelligencia használata.}
A dolgozat elkészítése során a szerző generatív mesterséges intelligencián alapuló eszközt (nagy nyelvi modellt) használt, kizárólag szövegszintű és szerkesztési célokra: a szöveg olvashatóságának és nyelvi minőségének javítására, a saját megfogalmazású szakaszok átfogalmazására és stilizálására, a fejezetek szerkezeti rendezésére, valamint a \LaTeX-formázás és a bibliográfiai hivatkozások formai egységesítésére. Az eszköz egyes szövegrészekhez kiindulási megfogalmazást is adott, amelyet a szerző minden esetben ellenőrzött és átdolgozott. A dolgozatban szereplő ábrák a szerző saját kódjából és adataiból készültek, nem pedig generatív eszközzel.

A dolgozat tudományos tartalma --- a HEMS-modell és a vizsgált optimalizáló algoritmusok megtervezése és implementálása, a kísérletek elvégzése, valamint az eredmények kiértékelése és értelmezése --- a szerző önálló munkája. A mesterséges intelligencia által adott kimeneteket a szerző ellenőrizte, és a dolgozat teljes tartalmáért, beleértve a hivatkozások helyességét, teljes felelősséget vállal. A felhasznált eszköz nem tekinthető a munka szerzőjének vagy társszerzőjének, és a bibliográfiában sem szerepel forrásként.
